% GENERATED FILE. Edit canonical lessons, then run npm run generate:books.
% canonical-source-hash: fnv1a64:dc67b65206d50a1d
% canonical-lessons: AR-C479-tazallaja, AR-C479-wahsh, AR-C479-kis, AR-C479-tabadala, AR-C479-hazama

\chapter{A day in the hills}
\label{ch:ar-a-day-in-the-hills}
\hlchaptermodality{479}

\begin{chapteropening}
\textbf{By the end of this chapter:} \emph{I can say to ski, a wild beast, a plastic bag, to exchange and to pack.}

Complete the last lesson of chapter 479: I can say to ski, a wild beast, a plastic bag, to exchange and to pack.
\end{chapteropening}

\section[\texorpdfstring{\ar{تزلج} (tazallaja)}{تزلج (tazallaja)}]{\ar{تزلج} (tazallaja) — to ski; to skate}

\label{lesson:AR-C479-tazallaja}



\begin{quote}
Before the new one: say the Arabic for fur, then the Arabic for to graze.
\end{quote}

\subsection*{You'll want to know: \ar{تزلج}}
\textbf{\ar{تزلج}} — \emph{tazallaja} — \textquotedblleft{}to ski; to skate\textquotedblright{}.

\begin{quote}
\ar{تزلجنا} \ar{على} \ar{الثلج} \ar{في} \ar{الشتاء}.
\end{quote}

\begin{quote}
\emph{tazallajnā ʿalā ath-thalj fī ash-shitāʾ.}
\end{quote}

\begin{quote}
We skied on the snow in winter.
\end{quote}

Present: \ar{يتزلج}, \emph{yatazallaj}.

\begin{cousinweb}
\textbf{\ar{تزلج}} (\emph{tazallaja}) is usually traced to the native Arabic root \emph{z-l-j}, \textquotedblleft{}to slip, to slide\textquotedblright{}. It is Form V.
\end{cousinweb}

\begin{grammarlens}[title={What you've built}]
One more word for getting around.
\end{grammarlens}

\subsection*{Your turn}
\begin{itemize}
\raggedright
  \item \emph{Say it:} \emph{tazallaja}
  \item \emph{Say it:} \emph{tazallaja}, once more
  \item \emph{Say it:} \emph{tazallajnā ʿalā ath-thalj fī ash-shitāʾ.} — We skied on the snow in winter.
  \item \emph{Recall:} say the Arabic for fur, then the Arabic for to graze, then say \emph{tazallaja} again
\end{itemize}

\begin{tcolorbox}[breakable,colback=teal!4,colframe=teal!35!black,title={Before you move on}]
What does \ar{تزلج} mean? (\textquotedblleft{}To ski; to skate\textquotedblright{}.) Say it once more.
\end{tcolorbox}
\section[\texorpdfstring{\ar{وحش} (waḥsh)}{وحش (waḥsh)}]{\ar{وحش} (waḥsh) — a wild beast; a monster}

\label{lesson:AR-C479-wahsh}



\begin{quote}
Before the new one: say the Arabic for to graze, then the Arabic for to ski.
\end{quote}

\subsection*{You'll want to know: \ar{وحش}}
\textbf{\ar{وحش}} — \emph{waḥsh} — \textquotedblleft{}a wild beast; a monster\textquotedblright{}.

\begin{quote}
\ar{خاف} \ar{الطفل} \ar{من} \ar{الوحش} \ar{في} \ar{القصة}.
\end{quote}

\begin{quote}
\emph{khāfa aṭ-ṭifl min al-waḥsh fī al-qiṣṣa.}
\end{quote}

\begin{quote}
The child was afraid of the beast in the story.
\end{quote}

Plural: \ar{وحوش}, \emph{wuḥūsh}.

\begin{cousinweb}
\textbf{\ar{وحش}} (\emph{waḥsh}) is built on the native Arabic root \emph{w-ḥ-sh}, \textquotedblleft{}wild, desolate\textquotedblright{} (Lane).
\end{cousinweb}

\begin{grammarlens}[title={What you've built}]
One more word for getting around.
\end{grammarlens}

\subsection*{Your turn}
\begin{itemize}
\raggedright
  \item \emph{Say it:} \emph{waḥsh}
  \item \emph{Say it:} \emph{waḥsh}, once more
  \item \emph{Say it:} \emph{khāfa aṭ-ṭifl min al-waḥsh fī al-qiṣṣa.} — The child was afraid of the beast in the story.
  \item \emph{Recall:} say the Arabic for to graze, then the Arabic for to ski, then say \emph{waḥsh} again
\end{itemize}

\begin{tcolorbox}[breakable,colback=teal!4,colframe=teal!35!black,title={Before you move on}]
What does \ar{وحش} mean? (\textquotedblleft{}A wild beast; a monster\textquotedblright{}.) Say it once more.
\end{tcolorbox}
\section[\texorpdfstring{\ar{كيس} (kīs)}{كيس (kīs)}]{\ar{كيس} (kīs) — a plastic bag; a sack}

\label{lesson:AR-C479-kis}



\begin{quote}
Before the new one: say the Arabic for to ski, then the Arabic for a wild beast.
\end{quote}

\subsection*{You'll want to know: \ar{كيس}}
\textbf{\ar{كيس}} — \emph{kīs} — \textquotedblleft{}a plastic bag; a sack\textquotedblright{}.

\begin{quote}
\ar{وضعنا} \ar{الطعام} \ar{في} \ar{كيس} \ar{كبير}.
\end{quote}

\begin{quote}
\emph{waḍaʿnā aṭ-ṭaʿām fī kīs kabīr.}
\end{quote}

\begin{quote}
We put the food in a big bag.
\end{quote}

Plural: \ar{أكياس}, \emph{akyās}.

\emph{ḥaqība} is a bag one carries; \emph{kīs} is a sack or a plastic bag.

\begin{cousinweb}
\textbf{\ar{كيس}} (\emph{kīs}) is given by Lane as a purse for money; the origin of the word is not settled.
\end{cousinweb}

\begin{grammarlens}[title={What you've built}]
One more word for getting around.
\end{grammarlens}

\subsection*{Your turn}
\begin{itemize}
\raggedright
  \item \emph{Say it:} \emph{kīs}
  \item \emph{Say it:} \emph{kīs}, once more
  \item \emph{Say it:} \emph{waḍaʿnā aṭ-ṭaʿām fī kīs kabīr.} — We put the food in a big bag.
  \item \emph{Recall:} say the Arabic for to ski, then the Arabic for a wild beast, then say \emph{kīs} again
\end{itemize}

\begin{tcolorbox}[breakable,colback=teal!4,colframe=teal!35!black,title={Before you move on}]
What does \ar{كيس} mean? (\textquotedblleft{}A plastic bag; a sack\textquotedblright{}.) Say it once more.
\end{tcolorbox}
\section[\texorpdfstring{\ar{تبادل} (tabādala)}{تبادل (tabādala)}]{\ar{تبادل} (tabādala) — to exchange (with each other)}

\label{lesson:AR-C479-tabadala}



\begin{quote}
Before the new one: say the Arabic for a wild beast, then the Arabic for a plastic bag.
\end{quote}

\subsection*{You'll want to know: \ar{تبادل}}
\textbf{\ar{تبادل}} — \emph{tabādala} — \textquotedblleft{}to exchange (with each other)\textquotedblright{}.

\begin{quote}
\ar{تبادلنا} \ar{العناوين} \ar{في} \ar{نهاية} \ar{الرحلة}.
\end{quote}

\begin{quote}
\emph{tabādalnā al-ʿanāwīn fī nihāyat ar-riḥla.}
\end{quote}

\begin{quote}
We exchanged addresses at the end of the trip.
\end{quote}

Present: \ar{يتبادل}, \emph{yatabādal}.

Form VI makes the action mutual: each gives to the other.

\begin{cousinweb}
\textbf{\ar{تبادل}} (\emph{tabādala}) is built on the native Arabic root \emph{b-d-l}, \textquotedblleft{}to change, to replace\textquotedblright{} (Lane). It is Form VI.
\end{cousinweb}

\begin{grammarlens}[title={What you've built}]
One more word for getting around.
\end{grammarlens}

\subsection*{Your turn}
\begin{itemize}
\raggedright
  \item \emph{Say it:} \emph{tabādala}
  \item \emph{Say it:} \emph{tabādala}, once more
  \item \emph{Say it:} \emph{tabādalnā al-ʿanāwīn fī nihāyat ar-riḥla.} — We exchanged addresses at the end of the trip.
  \item \emph{Recall:} say the Arabic for a wild beast, then the Arabic for a plastic bag, then say \emph{tabādala} again
\end{itemize}

\begin{tcolorbox}[breakable,colback=teal!4,colframe=teal!35!black,title={Before you move on}]
What does \ar{تبادل} mean? (\textquotedblleft{}To exchange (with each other)\textquotedblright{}.) Say it once more.
\end{tcolorbox}
\section[\texorpdfstring{\ar{حزم} (ḥazama)}{حزم (ḥazama)}]{\ar{حزم} (ḥazama) — to pack (luggage)}

\label{lesson:AR-C479-hazama}



\begin{quote}
Before the new one: say the Arabic for a plastic bag, then the Arabic for to exchange.
\end{quote}

\subsection*{You'll want to know: \ar{حزم}}
\textbf{\ar{حزم}} — \emph{ḥazama} — \textquotedblleft{}to pack (luggage)\textquotedblright{}.

\begin{quote}
\ar{حزمت} \ar{أمتعتي} \ar{قبل} \ar{الرحلة} \ar{بيوم}.
\end{quote}

\begin{quote}
\emph{ḥazamtu amtiʿatī qabla ar-riḥla bi-yawm.}
\end{quote}

\begin{quote}
I packed my luggage a day before the trip.
\end{quote}

Present: \ar{يحزم}, \emph{yaḥzim}.

\begin{cousinweb}
\textbf{\ar{حزم}} (\emph{ḥazama}) is built on the native Arabic root \emph{ḥ-z-m}, \textquotedblleft{}to bind, to tie up a bundle\textquotedblright{} (Lane). It shares its root with \emph{ḥizām}, \textquotedblleft{}a belt\textquotedblright{}.
\end{cousinweb}

\begin{grammarlens}[title={What you've built}]
One more word for getting around.
\end{grammarlens}

\subsection*{Your turn}
\begin{itemize}
\raggedright
  \item \emph{Say it:} \emph{ḥazama}
  \item \emph{Say it:} \emph{ḥazama}, once more
  \item \emph{Say it:} \emph{ḥazamtu amtiʿatī qabla ar-riḥla bi-yawm.} — I packed my luggage a day before the trip.
  \item \emph{Recall:} say the Arabic for a plastic bag, then the Arabic for to exchange, then say \emph{ḥazama} again
\end{itemize}

\begin{tcolorbox}[breakable,colback=teal!4,colframe=teal!35!black,title={Before you move on}]
What does \ar{حزم} mean? (\textquotedblleft{}To pack (luggage)\textquotedblright{}.) Say it once more.
\end{tcolorbox}
